\section{Las opciones en las que participas}
\begin{itemize}
  \item Artículo especializado para publicar en revista indizada
  \item Ensayo
  \item Memoria de experiencia laboral
  \item Reporte de aplicación de conocimientos
  \item Reporte de autoempleo profesional
  \item Reporte de residencia de investigación
  \item Tesina
  \item Tesis
\end{itemize}
\section{Lo que es igual en todas las opciones}
Las pantallas de esta sección se tomaron en la opción \emph{Tesis}; en las demás opciones son iguales.
\subsection{Firma de protocolo}
Abre el enlace, lee el documento con \boton{Leer el documento (PDF)} y pulsa \boton{Firmar}. El trámite avanza en cuanto firman todos los que deben hacerlo.
\captura{tesis_asesor_01_firma_de_protocolo}{Enlace personal del correo «Firma de protocolo»: lees el documento y pulsas «Firmar».}
\subsection{Vo. Bo. de la solicitud de revisión}
Abre el enlace, lee el documento con \boton{Leer el documento (PDF)} y pulsa \boton{Firmar}. El trámite avanza en cuanto firman todos los que deben hacerlo.
\captura{tesis_asesor_02_vo_bo_de_la_solicitud_de_revis}{Enlace personal del correo «Vo. Bo. de la solicitud de revisión»: lees el documento y pulsas «Firmar».}
\subsection{Firma del voto aprobatorio}
Abre el enlace, lee el documento con \boton{Leer el documento (PDF)} y pulsa \boton{Firmar}. El trámite avanza en cuanto firman todos los que deben hacerlo.
\captura{tesis_asesor_03_firma_del_voto_aprobatorio}{Enlace personal del correo «Firma del voto aprobatorio»: lees el documento y pulsas «Firmar».}
\subsection{Vo. Bo. de la solicitud de sustentación}
Abre el enlace, lee el documento con \boton{Leer el documento (PDF)} y pulsa \boton{Firmar}. El trámite avanza en cuanto firman todos los que deben hacerlo.
\captura{tesis_asesor_04_vo_bo_de_la_solicitud_de_suste}{Enlace personal del correo «Vo. Bo. de la solicitud de sustentación»: lees el documento y pulsas «Firmar».}
\section{Lo propio de cada opción}
\subsection{Artículo especializado para publicar en revista indizada}
Trabajo escrito de esta opción: Artículo original e inédito, enviado a una revista inscrita en índices validados. Firmas: protocolo, solicitud de revisión, voto aprobatorio y solicitud de sustentación.
\subsection{Ensayo}
Trabajo escrito de esta opción: Exposición crítica y propositiva de un tema, con postura intelectual propia. Firmas: protocolo, solicitud de revisión, voto aprobatorio y solicitud de sustentación.
\subsection{Memoria de experiencia laboral}
Trabajo escrito de esta opción: Recopilación y resultados de una experiencia profesional en un ámbito reconocido. Firmas: protocolo, solicitud de revisión, voto aprobatorio y solicitud de sustentación.
\subsection{Reporte de aplicación de conocimientos}
Trabajo escrito de esta opción: Propuesta de solución a un problema determinado, derivada de lo aprendido en la carrera. Firmas: protocolo, solicitud de revisión, voto aprobatorio y solicitud de sustentación.
\subsection{Reporte de autoempleo profesional}
Trabajo escrito de esta opción: Informe sobre la constitución y desarrollo de una unidad económica formal. Firmas: protocolo, solicitud de revisión, voto aprobatorio y solicitud de sustentación.
\subsection{Reporte de residencia de investigación}
Trabajo escrito de esta opción: Informe de actividades de apoyo a la investigación bajo la dirección de un investigador activo en el SNI. Firmas: protocolo, solicitud de revisión, voto aprobatorio y solicitud de sustentación.
\subsection{Tesina}
Trabajo escrito de esta opción: Investigación documental de carácter monográfico, con aportación al estado del arte. Firmas: protocolo, solicitud de revisión, voto aprobatorio y solicitud de sustentación.
\subsection{Tesis}
Trabajo escrito de esta opción: Investigación documental, experimental, empírica o teórica; individual o de hasta tres pasantes. Firmas: protocolo, solicitud de revisión, voto aprobatorio y solicitud de sustentación.
